\section{Introduction}
\label{sec:intro}

For integer $n,k\geq 2$ define the quantum (spatial) correlation set $C_{qs}(n,k)$ as the subset of $\R^{n^2k^2}$ that contains all tuples $(p_{abxy})$ representing families of bipartite distributions that can be locally generated in non-relativistic quantum mechanics. Formally, $(p_{abxy}) \in C_{qs}(n,k)$ if and only if there exist separable Hilbert spaces $\mH_\alice$ and $\mH_\bob$, for every $x\in\{1,\ldots,n\}$ (resp. $y\in\{1,\ldots,n\}$), a collection of projections $\{A^x_a\}_{a\in\{1,\ldots,k\}}$ on $\mH_\alice$ (resp. $\{B^y_b\}_{b\in\{1,\ldots,k\}}$ on $\mH_\bob$) that sum to identity, and a state (unit vector) $\psi \in \mH_\alice \otimes \mH_\bob$ such that 
\begin{equation}\label{eq:intro-qs}
  \forall x,y\in \{1, 2, \ldots, n\}\;,\quad\forall a,b\in\{1, 2, \ldots, k\}\;,
  \qquad p_{abxy} = \psi^* \big( A^x_a \otimes B^y_b\big) \psi\;.
\end{equation}
Note that due to the normalization conditions on $\psi$ and on $\{A^x_a\}$ and
$\{B^y_b\}$, for each $x,y$, $(p_{abxy})$ is a probability distribution on
$\{1, 2, \ldots, k\}^2$.
By taking direct sums it is easy to see that the set $C_{qs}(n,k)$ is convex. Let $C_{qa}(n,k)$ denote its closure (it is known that $C_{qs}(n,k) \neq C_{qa}(n,k)$, see~\cite{slofstra2019set}).


Our main result is that the family of sets $\{C_{qa}(n,k)\}_{n,k \in \N}$ is extraordinarily complex, in the following computational sense. For any $0<\eps<1$ define the $\eps$-\emph{weak membership problem} for $C_{qa}$ as the problem of deciding, given $n,k \in \N$ and a point $p=(p_{abxy}) \in \R^{n^2 k^2}$, whether $p$ lies in $C_{qa}(n,k)$ or is $\eps$-far from it in $\ell_1$ distance, promised that one is the case. Then we show that for any given $0<\eps<1$ the  $\eps$-{weak membership problem} for $C_{qa}$ cannot be solved by a Turing machine that halts with the correct answer on every input.

We show this by directly reducing the Halting problem to the weak membership problem for $C_{qa}$: we show that for all $0 < \eps < 1$ and any Turing machine $\cM$ one can efficiently compute integers $n,k\in \N$ and a linear functional $\ell_\cM$ on $\R^{n^2k^2}$ such that, whenever $\cM$ halts it holds that 
\begin{equation}\label{eq:intro-win-1}
\sup_{p\in C_{qa}(n,k)} \big|\ell_\cM(p)\big| \,=\,1\;,
\end{equation}
 whereas if $\cM$ does not halt then 
\begin{equation}\label{eq:intro-win-2}
\sup_{p\in C_{qa}(n,k)} \big|\ell_\cM(p)\big| \,\leq\, 1 - \eps\;.
\end{equation}
By standard results in convex optimization, this implies the aforementioned claim on the undecidability of the $\eps$-weak membership problem for $C_{qa}$ (for any $0<\eps<1$). 

Our result has interesting consequences for long-standing conjectures in quantum information theory and the theory of von Neumann algebras. Through a connection that follows from the work of Navascues, Pironio, and Acin~\cite{navascues2008convergent} the undecidability result implies a negative answer to Tsirelson's problem~\cite{Tsi06}. Let $C_{qc}(n,k)$ denote the set of quantum commuting correlations, which is the set of tuples $(p_{abxy})$ arising from operators $\{A^x_a\}$ and $\{B^y_b\}$ acting on a single Hilbert space $\mH$ and a state $\psi \in \mH$ such that
\begin{equation}\label{eq:intro-qc}
 \forall x,y\in \{1,\ldots,n\}\;,\;\forall a,b\in\{1,\ldots,k\}\;,\qquad p_{abxy} = \psi^* \big( A^x_a B^y_b\big) \psi \quad\text{and}\quad  \big[A^x_a,\,B^y_b\big]=0\;.
\end{equation}
Then \emph{Tsirelson's problem} asks if, for all $n,k$, the sets $C_{qa}(n,k)$ and $C_{qc}(n,k)$ are equal. 
Using results from~\cite{navascues2008convergent} we give integer $n,k$ and an explicit linear function $\ell$ on $\R^{n^2k^2}$ such that 
\[ \sup_{p\in C_{qc}(n,k)} \big|\ell(p)\big| \,=\,1\;,\qquad\text{but}\qquad \sup_{p\in C_{qa}(n,k)} \big|\ell(p)\big| \,\leq\,\frac{1}{2}\;,\]
which implies that $C_{qa}(n,k)\neq C_{qc}(n,k)$. 
By an implication of Fritz~\cite{fritz2012tsirelson} and Junge et al.~\cite{junge2011connes} we further obtain that Connes' Embedding Conjecture~\cite{connes1976classification} is false; in other words, there exist type II$_1$ von Neumann factors that do not embed in an ultrapower of the hyperfinite II$_1$ factor. We explain these connections in more detail in Section~\ref{sec:consequences} below.

Our approach to constructing such linear functionals on correlation sets goes through the theory of 
interactive proofs from complexity theory. 
To explain this connection we first review the concept of interactive proofs. The reader familiar with interactive proofs may skip the next section to arrive directly at a formal statement of our main complexity-theoretic result in Section~\ref{sec:result}. 

\subsection{Interactive proof systems}
An \emph{interactive proof system} is an abstraction that generalizes the familiar notion of \emph{proof}. Intuitively, given a formal statement $z$ (for example, ``this graph admits a proper $3$-coloring''), a proof $\pi$ for $z$ is information that enables one to check the validity of $z$ more efficiently than without access to the proof (in this example, $\pi$ could be an explicit assignment of colors to each vertex of the graph). 

Complexity theory formalizes the notion of proof in a way that emphasizes the role played by the verification procedure. To explain this, first recall that in complexity theory a \emph{language} $L$ is a subset of $\{0,1\}^*$, the set of all bit strings of any length, that intuitively represents all problem instances to which the answer should be ``yes''. For example, the language $L=\textsc{3-Coloring}$ contains all strings $z$ such that $z$ is the description (according to some pre-specified encoding scheme) of a $3$-colorable graph $G$. We say that a language $L$ admits efficiently verifiable proofs if there exists an algorithm $V$ (formally, a polynomial-time Turing machine) that satisfies the following two properties: (i) for any $z\in L$ there is a string $\pi$ such that $V(z,\pi)$ returns $1$ (we say that $V$ ``accepts''), and (ii) for any $z\notin L$ there is no string $\pi$ such that $V(z,\pi)$ accepts. Property (i) is generally referred to as the \emph{completeness} property, and (ii) is the \emph{soundness}. The set of all languages $L$ with both these completeness and soundness properties is denoted by the complexity class $\NP$.


Research in complexity and cryptography in the 1980s and 1990s led to a significant generalization of the notion of ``efficiently verifiable proof''. The first modification is to allow \emph{randomized} verification procedures by relaxing (i) and (ii) to \emph{high probability} statements: every $z\in L$ should have a proof $\pi$ that is accepted \emph{with probability at least $c$} (the completeness parameter), and for no $z\notin L$ should there be a proof $\pi$ that is accepted \emph{with probability larger than $s$} (the soundness parameter).
A common setting is to take $c=\frac{2}{3}$ and $s=\frac{1}{3}$; standard amplification techniques reveal that the exact values do not significantly affect the class of languages that admit such proofs, provided that they are chosen within reasonable bounds. 

The second modification is to allow \emph{interactive} verification. 
Informally, this means that instead of receiving a proof string $\pi$ in its entirety and making a decision based on it, the verification algorithm (called the ``verifier'') instead communicates with another algorithm called a ``prover'', and based on the communication decides whether $z\in L$. There are no restrictions on the computational power of the prover, whereas the verifier is required to run in polynomial time.\footnote{The reader may find the following mental model useful: in an interactive proof, an all-powerful prover is trying to convince a skeptical, but computationally limited, verifier that a string $z$ (known to both) lies in the set $L$, even when it may be that in fact $z \notin L$. By interactively interrogating the prover, the verifier can reject false claims, i.e.\ determine with high statistical confidence whether $z \in L$ or not. Importantly, the verifier is allowed to probabilistically and adaptively choose its messages to the prover. }




To understand how randomization and interaction can help for proof checking, consider the following example of an interactive proof for the language $\GNI$, which contains all pairs of graphs $(G_0,G_1)$ such that $G_0$ and $G_1$ are \emph{not} isomorphic.\footnote{Here and in the rest of the section, we implicitly assume that graphs and tuples of graphs have a canonical encoding as binary strings.}
It is not known if $\GNI\in \NP$, because it is not clear how to give an efficiently verifiable proof string that two graphs $G_0$ and $G_1$ are not isomorphic. (A proof of isomorphism is, of course, trivial: given a bijection from the vertices of $G_0$ to those of $G_1$ it is straightforward to verify that the bijection induces an isomorphism.) However, consider the following \emph{randomized}, \emph{interactive} verification procedure. Suppose the input to the verifier and prover is a pair of $n$-vertex graphs $(G_0,G_1)$ (if the graphs do not have the same number of vertices, they are trivially non-isomorphic and the verification procedure can automatically accept). The verifier first selects a uniformly random $b\in \{0,1\}$ and a uniformly random permutation $\sigma$ of $\{1,\ldots,n\}$ and sends the graph $H = \sigma( G_b)$ to the prover. The prover is then supposed to respond with a bit $b'\in \{0,1\}$; if $b'=b$  the verifier accepts and if $b'\neq b$ it rejects. 

Clearly, if $G_0$ and $G_1$ are not isomorphic then there exists a prover {strategy} to compute $b$ from $H$ with probability $1$: using its unlimited computational power, the prover can determine whether $H$ is isomorphic to $G_0$ or to $G_1$. However, if $G_0$ and $G_1$ \emph{are} isomorphic then the distribution of $H$ is uniform over the isomorphism class of $G_0$, which is the same as the isomorphism class of $G_1$, and the prover (despite having unlimited computational power) cannot distinguish between whether the verifier generated $H$ using $G_0$ or $G_1$. Thus the probability that \emph{any} prover can correctly guess $b'=b$ is exactly $\frac{1}{2}$. As a result, we have shown that the graph non-isomorphism problem has an interactive proof system with completeness $c=1$ and soundness $s=\frac{1}{2}$. Note how little ``information'' is communicated by the prover: a single bit!
The extreme succinctness of the ``proof'' comes from the fact that whether $G_0$ is isomorphic to $G_1$ determines whether a prover can reliably compute, given the data available to it (which is $G_0$,$G_1$, and $H$), the correct bit $b$.

We denote by $\IP$ the class of languages that admit randomized interactive proof systems such as the one just described. 
The class $\IP$ is easily seen to contain $\NP$, but it is thought to be a much larger class: one of the famous results of complexity theory is that $\IP$ is exactly the same as $\mathsf{PSPACE}$~\cite{lund1990algebraic,shamir1990ip}, the class of languages decidable by Turing machines using polynomial space.\footnote{The reason $\mathsf{PSPACE}$ is considered a ``difficult'' class of problems is because many computational problems believed to require super-polynomial or exponential time (such as $\textsc{3-Coloring}$ or deciding whether a quantified Boolean formula is true) can be solved using a polynomial amount of space.} Thus a polynomial-time verifier, when augmented with the ability to interrogate an all-powerful prover and use randomization, can solve computational problems that are (believed to be) vastly more difficult than those that can be checked using static, deterministic proofs (i.e. $\mathsf{NP}$ problems).

\paragraph{Multiprover interactive proofs.}
We now discuss a generalization of interactive proofs called \emph{multiprover interactive proofs}. Here, a polynomial-time verifier can interact with \emph{two} (or more) provers to decide whether a given instance $z$ is in a language $L$ or not. In this setting, after the verifier and all the provers receive the common input $z$, the provers are not allowed to communicate with each other, and the verifier ``cross-interrogates'' the provers in order to decide if $z \in L$. The provers may coordinate a joint strategy ahead of time, but once the protocol begins the provers can only interact with the verifier. As we will see, the extra condition that the provers cannot communicate with each other is a powerful constraint that can be leveraged by the verifier.





Consider the computational problem of deciding membership in a \emph{promise} language called $\GX$. 
A promise language $L$ is specified by two disjoint subsets $L_{yes},L_{no}\subseteq\{0,1\}^*$, and the task is to decide whether a given instance $z$ is in $L_{yes}$ or $L_{no}$, promised that $z \in L_{yes} \cup L_{no}$. In a proof system for a promise language, the completeness case consists of accepting with probability at least $c$ if $z \in L_{yes}$, and the soundness case consists of accepting with probability at most $s$ if $z \in L_{no}$. If $z \notin L_{yes} \cup L_{no}$, then there are no constraints on the behavior of the verifier.






The promise language $\GX$ is defined as follows: $\GX_{yes}$ (resp. $\GX_{no}$) is the set of all graphs $G$ with a \emph{cut} (i.e. a bipartition of the vertices) such that at least $90\%$ of edges cross the cut (resp. at most $60\%$ of edges cross the cut).\footnote{The specific numbers $90\%$ and $60\%$ are not too important; the only thing that really matters is that the first one is strictly less than $100\%$ and the second strictly larger than $50\%$, as otherwise the problem becomes much easier.} For simplicity, we also assume that all graphs in $\GX_{yes}\cup \GX_{no}$ are regular, i.e.\ the degree is a constant across all vertices in the graph. 

The $\GX$ problem clearly lies in $\NP$, since given a candidate cut it is easy to count the number of edges that cross it and verify that it is at least $90\%$ of the total number of edges. Observe that the length of the proof and the time required to verify it are linear in the size of the graph (the number of vertices and edges). \emph{Finding} the proof is of course much harder, but we are only concerned with the complexity of the verification procedure.

Now consider the following simple \emph{two-prover interactive proof system} for $\GX$. Given a graph $G$, the verification procedure first samples a uniformly random edge $e=\{u,v\}$ in $G$. It then sends a uniformly random $x\in\{u,v\}$ to the first prover, and a uniformly random $y\in \{u,v\}$ to the second prover. Each prover sees its respective question only and is expected to respond with a single bit, $a,b\in\{0,1\}$ respectively. The verification procedure accepts if and only if $a=b$ if $x=y$, and $a\neq b$ if $x\neq y$.

We claim that the verification procedure described in the preceding paragraph is a \emph{multiprover interactive proof system} for the language $\GX$, with completeness $c=0.95$ and soundness $s=0.9$, in the following sense. First, whenever $G\in \GX_{yes}$ then there is a successful strategy for the provers: specifically, the provers can fix an optimal bipartition and consistently answer ``$0$'' when asked about a vertex from one side of the partition, and ``$1$'' when asked about a vertex from the other side; assuming there exists a cut that is crossed by at least $90\%$ of the edges, this strategy succeeds with probability at least $\frac{1}{2} + \frac{1}{2}0.9$, where the first factor $\frac{1}{2}$ arises from the case when both provers are sent the same vertex, in which case they always succeed. 

Conversely, suppose given a strategy for the provers that is accepted with probability $p = \frac{1}{2} + \frac{1}{2}(1-\delta)$ when the verification procedure is executed on a (regular) $n$-vertex graph $G$. 
We then claim that $G$ has a cut crossed by at least a $1-2\delta$ fraction of all edges. To show this, we leverage the non-communication assumption on the provers. Since either prover's question is always a single vertex, their strategy can be represented by a function from the vertices of $G$ to answers in $\{0,1\}$. Any such function specifies a bipartition of $G$. While the provers'  bipartitions need not be identical, the fact that they succeed with high probability, for the case when they are sent the same vertex, implies that they must be consistent with high probability. Finally, the fact that they also succeed with high probability when sent opposite endpoints of a randomly chosen edge implies that either prover's bipartition must be cut by a large number of edges. Taking the contrapositive establishes the soundness property.

We denote by $\MIP$ the class of languages that have multiprover interactive proof systems such as the one described in the preceding paragraph. 
Note that, in comparison to the $\NP$ verification procedure for $\GX$ considered earlier, the interactive, two-prover verification is much more efficient in terms of the effort required for the verifier.
 Assuming the graph is provided in a convenient format,\footnote{For example, the graph can be specified via a circuit that takes as input an edge index --- using some arbitrary ordering --- and returns labels for the two endpoints of the edge.} it is possible to sample a random edge and verify the provers' answers in time and space that scales \emph{logarithmically} with the size of the graph. This exponential improvement in the efficiency of the verification procedure serves as the starting point for another celebrated result from complexity theory: $\mathsf{MIP}$ is exactly the same as the class $\mathsf{NEXP}$~\cite{babai1991non}, which are problems that admit \emph{exponential-time checkable proofs}.\footnote{An example of such a problem is the language $\textsc{Succinct-3-Coloring}$, which contains descriptions of polynomial-size circuits $C$ that specify a $3$-colorable graph $G_C$ on \emph{exponentially many} vertices.} The class $\NEXP$ contains $\PSPACE$, but is believed to be much larger; this suggests that the ability to interrogate more than one prover enables a polynomial-time verifier to verify much more complex statements.
 


\paragraph{Nonlocal games.}
In this paper we will only be concerned with multiprover interactive proof systems that consist of a single round of communication with two provers: the verifier first sends its questions to each of the provers, the provers respond with their answers, and the verifier decides whether to accept or reject. The class of problems that admit such interactive proofs is denoted $\MIP(2,1)$, and it is known that $\MIP=\MIP(2,1)$~\cite{feige1992two}. Such proof systems have a convenient reformulation using the language of \emph{nonlocal games}, that we now explain. 

In a nonlocal game, we say that a verifier interacts with multiple non-communicating \emph{players} (instead of provers --- there is no formal difference between the two terms). An $n$-question, $k$-answer nonlocal game $\game$ is specified by two procedures: a \emph{question sampling} procedure that samples a pair of questions $(x,y)\in \{1,\ldots,n\}^2$ for the players according to a distribution $\mu$ (known to the verifier and the players), and a \emph{decision} procedure that takes as input the players' questions and their respective answers $a,b\in\{1,\ldots,k\}$ and evaluates a predicate $D(x,y,a,b)\in \{0,1\}$ to determine the verifier's acceptance or rejection. In classical complexity theory, the main quantity associated with a nonlocal game $\game$ is its \emph{classical value}, which is the maximum success probability that two cooperating but non-communicating players have in the game. Formally, the classical value is defined as
\begin{equation}\label{eq:class-val}
 \val(\game) = \sup_{p\in C_c(n,k)} \sum_{x,y} \mu(x,y) \sum_{a,b} D(x,y,a,b) p_{abxy}\;,
\end{equation}
where the set $C_c(n,k)$ is the set of \emph{classical correlations}, which are tuples $(p_{abxy})$ such that there exists a set $\Lambda$ with probability measure $\nu$ and for every $\lambda \in \Lambda$ functions $A^\lambda,B^\lambda: \{1,2,\ldots,n\} \to \{1,2,\ldots,k\}$ such that
\[
	\forall x,y \in \{1,2,\ldots,n\}, \quad \forall a,b \in \{1,2,\ldots,k\}, \quad p_{abxy} = \Pr_{\lambda \sim \nu} (A^\lambda(x) = a \wedge B^\lambda(y) = b).
\]
This definition captures the intuitive notion that a classical strategy for the players is specified by (i) a distribution $\nu$ on $\Lambda$ that represents some probabilistic information shared by the players that is independent of the verifier's questions, and (ii) two functions $A^\lambda,B^\lambda$ that represent each players' ``local strategy'' for answering given their shared randomness $\lambda$ and question $x$ or $y$. \footnote{For the functional analyst we briefly note that if we define a tensor 
\[L  = \sum_{x,y,a,b} \mu(x,y)D(x,y,a,b) e_{xa}\otimes e_{yb} \in \R^{nk}\otimes \R^{nk}\]
then $\val(\game) = \|L\|_{\ell_1^n(\ell_\infty)^k \otimes_\epsilon \ell_1^n(\ell_\infty)^k}$, with $\otimes_\epsilon$ denoting the injective tensor norm of Banach spaces. (For more connections between interactive proofs, nonlocal games and tensor norms we refer to the survey~\cite{palazuelos2016survey}.)} Note that due to the shared randomness $\lambda$, the set $C_c(n,k)$ is a (closed) convex subset of $[0,1]^{n^2k^2}$. 

To make the connection with interactive proof systems, observe that the assertion that $L\in \MIP(2,1)$ precisely amounts to the specification of an efficient mapping\footnote{Here by ``efficient'' we mean that there should be a polynomial-time Turing machine that on input $z$ returns (i) a polynomial-size randomized circuit that samples from $\mu$, and (ii) a polynomial-size circuit that evaluates the predicate $D$.} from problem instances $z$ to games $\game_z$ such that whenever $z\in L$ then  $\val(\game_z) \geq \frac{2}{3}$, whereas if $z \notin L$ then $\val(\game_z) \leq \frac{1}{3}$. Thus the complexity of the optimization problem~\eqref{eq:class-val} captures the complexity of the decision problem $L$. The aforementioned characterization of $\MIP$ as the class $\NEXP$ by~\cite{babai1991non} shows that in general this optimization problem is very difficult: it is as hard as deciding any language in $\NEXP$. 



\subsection{Statement of result}
\label{sec:result}

We now introduce the main complexity class that is the focus of this paper: $\MIP^*$, the  ``entangled-prover'' analogue of the class $\MIP$ considered earlier. Informally the class $\MIP^*$, first introduced in~\cite{cleve2004consequences}, contains all languages that can be decided by a classical polynomial-time verifier interacting with multiple \emph{quantum} provers sharing \emph{entanglement}. We focus on the class $\MIP^*(2,1)$, which corresponds to the setting of one-round protocols with two  provers. Equivalently, a language $L$ is in $\MIP^*(2,1)$ if and only if there is an efficient mapping from instances $z\in\{0,1\}^*$ to nonlocal games $\game_z$ such that if $z \in L$, then $\val^*(\game_z) \geq 2/3$ and otherwise $\val^*(\game_z) \leq 1/3$. Here, for an $n$-question, $k$-answer game $\game$, we let $\val^*(\game)$ denote its \emph{entangled value}, which is defined as
\begin{equation}\label{eq:ent-val}
 \val^*(\game) = \sup_{p\in C_{q}(n,k)} \sum_{x,y} \mu(x,y) \sum_{a,b} D(x,y,a,b) p_{abxy}\;,
\end{equation}
where the set $C_{q}(n,k)$ is the set of all finite-dimensional quantum correlations, i.e.\ correlations of the form~\eqref{eq:intro-qs} where $\mH$ is restricted to be finite-dimensional. Although the sets $C_q(n,k)$ and $C_{qs}(n,k)$ in general are distinct~\cite{coladangelo2018unconditional}, it is an easy exercise to verify that they have the same closure $C_{qa}(n,k)$, and therefore the supremum in~\eqref{eq:ent-val} can equivalently be taken over $C_q(n,k)$, $C_{qs}(n,k)$ or $C_{qa}(n,k)$. We use $C_q$ for convenience in the analysis.


Since $C_c(n,k) \subseteq C_{q}(n,k)$ for all $n,k$, we have that $\val(\game) \leq \val^*(\game)$; in other words, quantum spatial strategies can perform at least as well as classical strategies in a nonlocal game. 



The consideration of quantum strategies and the set $C_{qs}(n,k)$ for the definition of $\MIP^*$ is motivated by a long line of works in the foundations of quantum mechanics around the topic of \emph{Bell inequalities}, that are linear functionals which separate the sets $C_{c}(n,k)$ and $C_{qs}(n,k)$. The simplest such functional is the \emph{CHSH inequality}~\cite{clauser1969proposed}, that shows $C_c(2,2) \subsetneq C_{qs}(2,2)$. The CHSH inequality can be reformulated as a game $\game$ such that $\val^*(\game) > \val(\game)$. This game is very simple: it is defined by setting $\mu(x,y)=\frac{1}{4}$ for all $x,y\in\{0,1\}$ and $D(x,y,a,b)=1$ if and only if $a\oplus b = x\wedge y$. It can be shown that $\val(\game) = \frac{3}{4}$ and $\val^*(\game) = \frac{1}{2} + \frac{1}{2\sqrt{2}} > \frac{3}{4}$. The study of Bell inequalities is a large area of research not only in foundations, where they are a tool to study the nonlocal properties of entanglement, but also in quantum cryptography, where they form the basis for cryptographic protocols for e.g. quantum key distribution~\cite{ekert1991quantum}. 

The introduction of entanglement in the setting of interactive proofs has interesting consequences for complexity theory; indeed it is not \emph{a priori} clear how the class $\MIP^*$ compares to $\MIP$. Take a language $L\in \MIP(2,1)$, and let $z$ be an instance. Then the associated game $\game_z$ is such that $\val(\game_z)\geq\frac{2}{3}$ if $z\in L$, and $\val(\game_z)\leq\frac{1}{3}$ otherwise. The fact that in general $\val^*(\game_z) \geq \val(\game_z)$ (and that as demonstrated by the CHSH game inequality can be strict) cuts both ways. On the one hand, the soundness property can be affected, so that instances $z\notin L$ could have $\val^*(\game_z) = 1$, meaning that we would not be able to establish that $L \in \MIP^*$.
On the other hand, a language $L \in \MIP^*(2,1)$ may not necessarily be in $\MIP$, because for $z\in L$ the fact that $\val^*(\game_z)\geq\frac{2}{3}$ does not automatically imply $\val(\game_z)> \frac{1}{3}$ (in other words, the game $\game_z$ may require the players to use a quantum strategy in order to win with probability greater than $1/3$). Just as the complexity of the class $\MIP$ is characterized by the complexity of approximating the classical value of nonlocal games (the optimization problem in~\eqref{eq:class-val}), the complexity of $\MIP^*$ is intimately related to the complexity of approximating the entangled value of games (the optimization problem in~\eqref{eq:ent-val}).


In~\cite{ito2012multi} the first non-trivial lower bound on $\MIP^*$ was shown, establishing that $\MIP = \NEXP \subseteq \MIP^*$. (Earlier results~\cite{kempe2011entangled,ito2009oracularization} gave more limited hardness results, for approximating the entangled value up to inverse polynomial precision.) This was proved by arguing that for the specific games constructed by~\cite{babai1991non} that show $\NEXP \subseteq \MIP$, the classical and entangled values are approximately the same. In other words, the classical soundness and completeness properties of the proof system of~\cite{babai1991non} are maintained in the presence of shared entanglement between the provers. 
Following~\cite{ito2012multi} a sequence of works~\cite{vidick2016three,ji2016classical,natarajan2018two,ji2017compression,natarajan2018low,fitzsimons2018quantum} established progressively stronger lower bounds on the complexity of approximating the entangled value of nonlocal games, culminating in~\cite{NW19} which showed that approximating the entangled value is at least as hard as $\NEEXP$, the collection of languages decidable in non-deterministic \emph{doubly exponential} time. This proves that $\NEEXP \subseteq \MIP^*$, and since it is known that $\NEXP\subsetneq \NEEXP$ it follows that $\MIP \neq \MIP^*$.




In contrast to these increasingly strong lower bounds the only upper bound known on $\MIP^*$ is the trivial inclusion $\MIP^* \subseteq \RE$, the class of recursively enumerable languages, i.e. languages $L$ such that there exists a Turing machine $\cM$ such that $x\in L$ if and only if $\cM$ halts and accepts on input $x$. This inclusion follows since the supremum in~\eqref{eq:ent-val} can be approximated from below by performing an exhaustive search in increasing dimension and with increasing accuracy. We note that, in addition to containing all decidable languages, this class also contains undecidable problems such as the Halting problem, which is to decide whether a given Turing machine eventually halts.

Our main result is a proof of the reverse inclusion: $\RE \subseteq \MIP^*$. Combined with the preceding observation it follows that 
\[ \MIP^* \,=\, \RE \;,\]
which is a full characterization of the power of entangled-prover interactive proofs. In particular for any $0 < \eps < 1$, it is an undecidable problem to determine whether a given nonlocal game has entangled value $1$ or at most $1 - \eps$ (promised that one is the case).

\paragraph{Proof summary.}
The proof of the inclusion $\RE \subseteq \MIP^*$ is obtained by designing an entangled-prover interactive proof for the Halting problem, which is complete for the class $\RE$. Specifically, we design an efficient transformation that maps any Turing machine $\cal{M}$ to a 
nonlocal game $\game_{\cal{M}}$ such that, if $\cal{M}$ halts (when run on an empty input tape) then there is a quantum strategy for the provers that succeeds with probability $1$ in $\game_{\cal{M}}$ (i.e.\ $\val^*(\game_{\cal{M}})=1$), whereas if $\cal{M}$ does not halt then no quantum strategy can succeed with probability larger than $\frac{1}{2}$ in the game (i.e.\ $\val^*(\game_{\cal{M}})\leq \frac{1}{2}$). Furthermore, the game $\game_\cal{M}$ has the property that whenever $\val^*(\game_\cal{M})=1$ then this fact is witnessed by a \emph{synchronous} strategy, i.e.\ a strategy where the players always give the same answer when simultaneously asked the same question. Synchronous strategies, or correlations, were first introduced in~\cite{paulsen2016estimating} and have played an important role in approaches to CEP based on quantum information and the study of nonlocal games; see e.g.~\cite{kim2018synchronous} and references therein. 
In the paper we use a terminology of \emph{projective, consistent and commuting} (PCC) strategies (see Definition~\ref{def:spcc} in Section~\ref{sec:strat}), a notion which implies the notion of being synchronous. 



A very rough sketch of this construction is as follows (we give a detailed overview in Section~\ref{sec:overview}). Given an infinite family of games $\{ \game_n \}_{n \in \N}$, we say that the family is \emph{uniformly generated} if there is a polynomial-time Turing machine that on input $n$ returns a description of the game $\game_n$. Given a game $\game$ and $p\in[0,1]$ let $\Ent(\game,p)$ denote the minimum local dimension of an entangled state shared by the players in order for them to succeed in $\game$ with probability at least $p$. 

We proceed in two steps. First, we design a \emph{compression} procedure for a specific class of nonlocal games that we call \emph{normal form}. Given as input a uniformly generated family $\{ \game_{n} \}_{n \in \N}$ of normal form games, the compression procedure returns another uniformly generated family $\{\game'_{n}\}_{n \in \N}$ of normal form games with the following properties: (i) for all $n$, if $\val^*(\game_{2^n})=1$ then $\val^*(\game'_n)=1$, and (ii) for all $n$, if $\val^*(\game_{2^n})\leq \frac{1}{2}$ then $\val^*(\game'_n)\leq \frac{1}{2}$ and moreover
\[ \Ent(\game'_{n},\frac{1}{2})\,\geq\, \max \Big\{ \Ent\Big(\game_{2^n},\frac{1}{2}\Big), \,2^{2^{\Omega(n)}} \Big\}\;.\] 

The construction of this compression procedure is our main contribution. Informally, it combines the recursive compression technique developed in~\cite{ji2017compression,fitzsimons2018quantum} with the so-called ``introspection'' technique of~\cite{NW19} that was used to prove $\NEEXP \subseteq \MIP^*$. The introspection technique itself relies heavily on the quantum low individual degree test of~\cite{natarajan2018low,ML20}\tnote{added ref. to ML doc here} to robustly self-test certain distributions that arise from constructions of \emph{classical} probabilistically checkable proofs. The quantum low-degree test and the introspection technique allow us to avoid the shrinking gap limitation of the results from~\cite{fitzsimons2018quantum}. 




In the second step, we use the compression procedure in an iterated fashion to construct an interactive proof system for the Halting problem. Fix a Turing machine $\cM$ and consider the following family  of nonlocal games $\{\game_{\cM,n}^{(0)}\}_{n\in\N}$: for all $n \in \N$, if $\cM$ halts in at most $n$ steps (when run on an empty input tape), then $\val^*(\game_{\cM,n}^{(0)})=1$, and otherwise  $\val^*(\game_{\cM,n}^{(0)}) \leq \frac{1}{2}$.\footnote{There is nothing special about the choice of $\frac{1}{2}$; this can be set to any constant that is less than $1$.} 

Constructing such a family of games is trivial; furthermore, they can be made in the  ``normal form'' required by the compression procedure. However, consider applying the compression procedure to obtain a family of normal form games $\{ \game_{\cM,n}^{(1)} \}_{n\in\N}$. Then for all $n \in \N$, it holds that if $\cM$ halts in at most $2^n$ steps then $\val^*(\game_{\cM,n}^{(1)})=1$, and otherwise $\val^*(\game_{\cM,n}^{(1)})\leq \frac{1}{2}$, and furthermore any strategy that achieves a value of at least $\frac{1}{2}$ requires an entangled state of dimension at least $2^{2^{\Omega(n)}}$.



Intuitively, one would expect that iterating this procedure and ``taking the limit'' gives a family of games $\{ \game_{\cM,n}^{(\infty)} \}_{n\in\N}$ such that if $\cM$ halts then $\val^*(\game_{\cM,n}^{(\infty)})=1$ for all $n \in \N$, whereas if $\cM$ does not halt then no finite-dimensional strategy can succeed with probability larger than $\frac{1}{2}$ in $\game_{\cM,n}^{(\infty)}$, for all $n \in \N$; in particular $\val^*(\game_{\cM,n}^{(\infty)})\leq \frac{1}{2}$. 
Formally, we do not take such a limit but instead define directly the family of games $\{ \game_{\cM,n}^{(\infty)} \}_{n\in\N}$ as a \emph{fixed point} of the Turing machine that implements the compression procedure. The game $\game_\cM$ can then be taken as $\game_{\cM,1}^{(\infty)}$. We describe this in more detail in Section~\ref{sec:overview}. 


	
	
	
	





	

\subsection{Consequences}
\label{sec:consequences}

	Our result is motivated by a connection with Tsirelson's problem from quantum information theory, itself related to Connes' Embedding Conjecture in the theory of von Neumann algebras~\cite{connes1976classification}. 
In a celebrated sequence of papers, Tsirelson~\cite{tsirelson1993some} initiated the systematic study of quantum correlation sets. Recall the definition of the set of quantum spatial correlations 
\begin{equation}
\label{eq:spatial}
	C_{qs}(n,k) = \big\{ (p_{abxy} ) \mid p_{abxy} = \bra{\psi} A^x_a \otimes B^y_b \ket{\psi}, \; \ket{\psi} \in  \mH_\alice \otimes \mH_\bob,\,\;\forall xy,\, \{A^x_a\}_a,\{B^y_b\}_b\text{ POVM}\big\}\;,
\end{equation}
where here $\ket{\psi}$ ranges over all unit norm vectors $\ket{\psi} \in \mH_\alice \otimes \mH_\bob$ with $\mH_\alice$ and $\mH_\bob$ arbitrary (separable) Hilbert spaces, and a POVM is defined as a collection of positive semidefinite operators that sum to identity. (From now on we use the Dirac ket notation $\ket{\psi}$ for states.) Recall the closure $C_{qa}(n,k)$ of $C_{qs}(n,k)$.

Tsirelson observed that there is a natural alternative definition to the quantum
spatial correlation set, called the \emph{quantum commuting correlation set} and
defined as
\begin{equation}
  \label{eq:commuting}
  C_{qc}(n,k) = \bigl\{ (p_{abxy}) \mid p_{abxy} = \bra{\psi} A^x_a\,  B^y_b
  \ket{\psi} \bigr\}\;,
\end{equation}
where $\ket{\psi}\in \mH$ is a quantum state, $\{A^x_a\}$ and $\{B^y_b\}$ are
POVMs for all $x, y$, and $[A^x_a, B^y_b] = 0$ for all $a, b, x, y$.
Note the key difference with spatial correlations is that
in~\eqref{eq:commuting} all operators act on the same (separable) Hilbert space.
The requirement that operators associated with different inputs (questions)
$x,y$ commute is arguably a minimal requirement within the context of quantum
mechanics for there to not exist any causal connection between outputs (answers)
$a,b$ obtained in response to the respective input.

The set $C_{qc}(n,k)$ is closed and convex, and it is easy to see that $C_{qa}(n,k)\subseteq C_{qc}(n,k)$ for all $n,k\geq 1$. When Tsirelson initially introduced these sets he claimed that equality holds. However, it was later pointed out that this is not obviously true. The question of equality between $C_{qc}$ and $C_{qa}$ (for all $n,k$) is now known as \emph{Tsirelson's problem}~\cite{Tsi06}. Let $	C_q(n,k) $ denote the same as $C_{qs}(n,k)$ except that both $\mH_\alice$ and $\mH_{\bob}$ in~\eqref{eq:spatial} are restricted to finite-dimensional spaces. Then more generally one can consider the following chain of inclusions
\begin{equation}
\label{eq:corr-sets-chain}
	C_q(n,k) \subseteq C_{qs}(n,k) \subseteq C_{qa}(n,k) \subseteq C_{qc}(n,k)\;,
\end{equation}
for all $n,k \in \N$, and ask which (if any) of these inclusions are strict. We let $C_q,C_{qs},C_{qa},C_{qc}$ denote the union of $C_q(n,k),C_{qs}(n,k),C_{qa}(n,k),C_{qc}(n,k)$, respectively, over all integers $n,k \in \N$. 
In a breakthrough work, Slofstra established the first separation between these four correlation sets by proving that $C_{qs} \neq C_{qc}$~\cite{slofstra2019tsirelson}; he later proved the stronger statement that $C_{qs} \neq C_{qa}$~\cite{slofstra2019set}. As a consequence of the technique used to demonstrate the separation Slofstra also obtains the complexity-theoretic statement that the problem of determining whether an element $p$ lies in $C_{qc}$, even promised that if it does, then it also lies in $C_{qa}$, is undecidable. Interestingly, this is shown by reduction from the \emph{complement} of the halting problem; for our result we reduce from the halting problem (see Section~\ref{sec:open} for further discussion of this point). 
Since his work, simpler proofs of Slofstra's results have been found~\cite{dykema2019non,musat2018non,coladangelo2019two}. In~\cite{coladangelo2018unconditional}, Coladangelo and Stark showed that $C_{q} \neq C_{qs}$ by exhibiting a $5$-input, $3$-output correlation that can be attained using infinite-dimensional spatial strategies (i.e.\ infinite-dimensional Hilbert spaces, a state and POVMs satisfying~\eqref{eq:spatial}) but cannot be attained via finite-dimensional strategies.

As already noted in~\cite{fritz2014can} (and further elaborated on by~\cite{fitzsimons2018quantum}), the undecidability of $\MIP^*(2,1)$ implies the separation $C_{qa} \neq C_{qc}$.\footnote{Technically~\cite{fritz2014can} make the observation for the commuting-prover analogue $\MIPco(2,1)$, discussed further in Section~\ref{sec:open}, but the reasoning is the same.} This follows from the observation that if $C_{qa} = C_{qc}$, then there exists an algorithm that can correctly determine if a nonlocal game $\game$ satisfies $\val^*(\game)=1$ or $\val^*(\game)\leq  \frac{1}{2}$ and always halts: this algorithm interleaves a hierarchy of semidefinite programs providing outer approximations to the set $C_{qc}$~\cite{navascues2008convergent,doherty2008quantum} with a simple exhaustive search procedure providing inner approximations to $C_{q}$. Our result that $\RE \subseteq \MIP^*(2,1)$ implies that no such algorithm exists, thus resolving Tsirelson's problem in the negative.

We furthermore exhibit an explicit nonlocal game $\game$ such that $\val^*(\game) < \valco(\game) = 1$, where $\valco(\game)$ is defined as $\val^*(\game)$ except that the supremum is taken over the set $C_{qc}(n,k)$ in~\eqref{eq:commuting}. This in turn yields an explicit correlation that is in the set $C_{qc}$ but not in $C_{qa}$. 
This game closely resembles the game $\game_\machine$ described in the sketch of the proof that $\RE \subseteq \MIP^*$, where $\machine$ is the Turing machine that runs the hierarchy of semidefinite programs on the game $\game_\machine$ and halts if it certifies that $\valco(\game_\machine) < 1$. It is in principle possible to determine an upper bound on the parameters $n,k$ for our separating correlation from the proof. While we do not provide such a bound, there is no step in the proof that requires it to be astronomical; e.g.\ we believe (without proof) that $10^{20}$ is a clear upper bound.

\paragraph{Connes' Embedding Conjecture.} 
Connes' Embedding Conjecture (CEC)~\cite{connes1976classification} is a conjecture in the theory of von Neumann algebras. Briefly, CEC posits that every type II$_1$ von Neumann factor embeds into an ultrapower of the hyperfinite II$_1$ factor. We refer to~\cite{ozawa2013connes} for a precise formulation of the conjecture and connections to other conjectures in operator algebras, such as  Kirchberg's QWEP conjecture. In independent work Fritz~\cite{fritz2012tsirelson} and Junge et al.~\cite{junge2011connes} showed that a positive answer to CEC would imply a positive resolution of Tsirelson's problem, i.e.\ that $C_{qa}(n,k)=C_{qc}(n,k)$ for all $n,k$. (This was later promoted to an equivalence by Ozawa~\cite{ozawa2013connes}.) Since our result disproves this equality for some $n,k$ it also implies that CEC does not hold. In work that appeared subsequently to the first announcement of our results, Goldbring and Hart~\cite{goldbring2016computability,goldbring2020universal} show using arguments from continuous logic that the uncomputability of approximating $\val^*(\game)$ refutes the CEC. Interestingly, their argument uses general elementary considerations from logic and completely bypasses Tsirelson's problem and its equivalence with Kirchberg's QWEP conjecture. We note that using the constructive aspect of our result it may be possible to give an explicit description of a factor that does not embed into an ultrapower of the hyperfinite II$_1$ factor, but we do not give such a construction. 

\paragraph{Entanglement tests.} As a step towards showing our result for any integer $n\geq 1$ we construct a game $\game_n$, with question and answer length polynomial in the size of the smallest Turing machine $\cal{M}_n$ that halts (on the empty tape) in exactly $n$ steps (i.e.\ the \emph{Kolmogorov complexity} of $n$), such that $\val^*(\game_n)=1$ yet any quantum strategy that succeeds in $\game_n$ with probability larger than $\frac{1}{2}$ must use an entangled state whose Schmidt rank is at least $2^{\Omega(n)}$. This is by far the most efficient entanglement test that we are aware of.

\paragraph{Prover and round reduction for $\MIP^*$ protocols.} Let $\MIP^*(k,r)$ denote the collection of languages decidable by $\MIP^*$ protocols with $k\geq 2$ provers and $r$ rounds. Prior to our work it was known how to perform \emph{round reduction} for $\MIP^*$ protocols, at the cost of adding provers; it was shown by~\cite{ji2017compression,fitzsimons2018quantum} that $\MIP^*(k,r) \subseteq \MIP^*(k + 15,1)$ for all $k, r$. However, it was an open question whether the complexity of the class $\MIP^*$ increases if we add more provers.
Our main complexity-theoretic result implies that $\MIP^* = \MIP^*(2,1)$. This follows from the following chain of inclusions: for all polynomially-bounded functions $k,r$,
 \[
 	\MIP^*(2,1) \subseteq \MIP^*(k,r) \subseteq \RE \subseteq \MIP^*(2,1)\;.
\]
The first inclusion follows since the verifier in an $\MIP^*$ protocol can always ignore extra provers and rounds; the second inclusion follows from a simple exhaustive-search procedure that enumerates over strategies for a given $\MIP^*(k,r)$ protocol; the third result is proven in this paper.\footnote{In fact, we note that the second term $\MIP^*(k,r)$ can be replaced by $\QMIP^*(k,r)$, which is the analogous class with a quantum verifier and quantum messages, since the first inclusion is trivial and the second remains true. As a result, we obtain that $\QMIP^* = \MIP^*(2,1)$ as well.}

However, this method of reducing provers and rounds in a given $\MIP^*$ protocol is indirect; it involves first converting a given $\MIP^*$ protocol into a Turing machine that accepts if and only if the $\MIP^*$ protocol has value larger than $\frac{1}{2}$, and then constructing an $\MIP^*(2,1)$ protocol to decide whether the Turing machine halts. In particular this transformation does not generally preserve the complexity of the provers and verifier in the original protocols. We leave it as an open question to find a more direct method for reducing the number of provers in an $\MIP^*$ protocol.



\subsection{Open questions}
\label{sec:open}

We mention several questions left open by our work.

\paragraph{Explicit constructions of counter-examples to Connes' Embedding Conjecture.} We provide an explicit counter-example to Tsirelson's problem in the form of a game whose entangled value differs from its commuting-operator value. Through the aforementioned connection with Connes' embedding conjecture~\cite{fritz2012tsirelson,junge2011connes,ozawa2013connes}, the counter-example may lead to the construction of interesting objects in other areas of mathematics. A first question is whether it can lead to an explicit description of a type II$_1$ factor that does not satisfy the Connes embedding property. Such a construction could be obtained along the lines of~\cite{kim2018synchronous}, using the fact that our game $\game$ such that $\val^*(\game) < \valco(\game) = 1$ has the property of being \emph{synchronous}, i.e.\ perfect strategies in the game are required to return the same answer when both parties are provided the same question. 

Going further, one may ask if the example can eventually lead to a construction of a group that is not sofic, or even not hyperlinear (see e.g.~\cite{capraro2015introduction} for the connection). 

\paragraph{The complexity of variants of $\MIP^*$.} Our result characterizes the complexity class $\MIP^*$ as the set of recursively enumerable languages. One can also consider the complexity class $\MIPco$, which stands for \emph{multiprover interactive proofs in the commuting-operator model}. For the sake of the discussion we consider only two-prover one-round protocols; a language $L$ is in $\MIPco(2,1)$ if there exists an efficient reduction that maps $z \in \{0,1\}^*$ to a nonlocal game $\game_z$ such that if $z \in L$ then $\valco(\game_z) \geq \frac{2}{3}$, and otherwise $\valco(\game_x) \leq \frac{1}{3}$.

The semidefinite programming hierarchy of~\cite{navascues2008convergent,doherty2008quantum} can be used to show that $\MIPco(2,1)$ is contained in the \emph{complement} of $\RE$, denoted as $\coRE$: to certify that $z \notin L$ it suffices to run the hierarchy until it obtains a certificate that $\valco(\game_z) < \frac{2}{3}$. Since it is known that $\RE \neq \coRE$,\footnote{$\RE \neq \coRE$ follows from the fact that $\RE \cap \coRE$ is the set of decidable languages and $\RE$ contains undecidable languages.} this implies that $\MIP^*(2,1) \neq \MIPco(2,1)$.

It is thus plausible that $\MIPco = \coRE$,\footnote{We note that the ``co'' modifier on both sides of the equation $\MIPco = \coRE$ refer to different things!} which would provide a very pleasing ``dual'' complexity characterization to $\MIP^* = \RE$. One possible route to proving this would be to adapt our gap-preserving compression framework to the commuting-operator setting by showing that each of the steps (question reduction, answer reduction, and parallel repetition) remain sound against commuting-operator strategies. Using the connection established in~\cite{fritz2014can}, this would imply that the operator norm over the maximal $C^*$ algebra $C^*(F_2 * F_2)$, where $F_2$ is the free group on two elements, is uncomputable. 

Another interesting open question concerns the \emph{zero gap} variants of $\MIP^*$ and $\MIPco$, which we denote by $\MIP^*_0$ and $\MIPco_0$, respectively. These classes capture the complexity of deciding whether a nonlocal game $\game$ has entangled value (or commuting-operator value respectively) \emph{exactly} equal to $1$. In~\cite{slofstra2019set}, Slofstra shows that there is an efficient reduction from Turing machines $\cal{M}$ to nonlocal games $\game_\machine$ such that $\machine$ does not halt if and only if $\val^*(\game_\machine) = \valco(\game_\machine) = 1$. This implies that $\coRE = \MIPco_0(2,1)$ and furthermore $\coRE \subseteq \MIP^*_0$. However, since $\RE \subseteq \MIP^*(2,1) \subseteq \MIP^*_0(2,1)$, this implies that $\MIP^*_0(2,1)$ is \emph{strictly} bigger than both $\RE$ and $\coRE$. In fact, it was shown by Mousavi, et al.~\cite{mousavi2020complexity} that the class $\MIP^*_0(3,1)$ (the complexity class corresponding to determining whether $\val^*(\game) = 1$ for \emph{three-player} nonlocal games) is equal to $\Pi_2^0$, which is the set of all languages $L$ such that $x \in L$ if and only if $\forall\, y, \, \exists z \, R(x,y,z) = 1$ for a computable function $R$ that depends on $L$.\footnote{The class $\Pi_2^0$ is also characterized as being part of the second level of the \emph{arithmetical hierarchy} from computability theory, where $\RE = \Sigma_1^0$ and $\coRE = \Pi_1^0$ form the first level.} Thus it is plausible that the complexity landscape of nonlocal games looks like the following: $\MIP^* = \RE$, $\MIP^*_0 = \Pi_2^0$, and $\MIPco = \MIPco_0 = \coRE$. Such statements about the complexity of $\MIP^*$ versus $\MIPco$, in both the gapped and zero-gap cases, may reveal additional insights into the difference between the tensor product and commuting-operator models of correlations.



\paragraph{Acknowledgments.} We thank Matthew Coudron, William Slofstra and Jalex Stark for enlightening discussions regarding possible consequences of our work. We thank William Slofstra and Jalex Stark for suggestions regarding explicit separations between $C_{qa}$ and $C_{qc}$. We thank Peter Burton, William Slofstra and Jalex Stark for comments on a previous version. We thank Lewis Bowen and Mikael de la Salle for pointing out typos and minor errors in a previous version.

Zhengfeng Ji is supported by Australian Research Council (DP200100950). Anand Natarajan is supported by IQIM, an NSF Physics Frontiers Center (NSF Grant PHY-1733907). Thomas Vidick is supported by NSF CAREER Grant CCF-1553477, AFOSR YIP award number FA9550-16-1-0495, a CIFAR Azrieli Global Scholar award, MURI Grant FA9550-18-1-0161 and the IQIM, an NSF Physics Frontiers Center (NSF Grant PHY-1125565) with support of the Gordon and Betty Moore Foundation (GBMF-12500028). Henry Yuen is supported by NSERC Discovery Grant 2019-06636.
Part of this work was done while John Wright was at the Massachusetts Institute of Technology.
He is supported by IQIM, an NSF Physics Frontiers Center (NSF Grant PHY-1733907), and by ARO contract W911NF-17-1-0433.



 
